\section{Discussion \& Limitations}
\label{sec:discussion}

The current Block makes computational operations and their invocation explicit
within a trainable continuous system. Its contextual sequence development run
shows that terminal supervision can train useful nonlinear state updates and
calling behavior. The same-size static-caller model is also strong. Because
that control still gives the primitives and memory full contextual inputs,
these observations do not isolate a benefit of dynamic calling. The useful
next mechanistic question is which learned primitive functions, memory reads
and invocation decisions account for the acquired computation, and whether
those roles remain reusable when task demands change.

Parameter-matched character development makes the comparison more specific.
After the same 300 updates, the 41,593-parameter dynamic Block has sampled
perplexities 12.458, 12.468 and12.496 at contexts 32, 64 and 128, respectively.
The static variant has 12.383, 12.452 and12.436. The 43,300-parameter Transformer
has 12.803, 49.837 and48.409, while the 41,548-parameter GRU has approximately
10.090 across these contexts. Context 64 is the designated unseen-length split
and context 128 the long split. These are 512 paired sampled positions per split,
not whole-corpus or word-level perplexities. They show a stable contextual
behavior for the tested Block but also a stronger recurrent control. The
Transformer's long-prefix result characterizes this configuration and budget,
not the whole architecture family. Matching parameters does not match
training work, forward cost or transient memory.

Initial real-input diagnostics expose a further gap. With the first ten
official classes, at most 128 training examples per class and 300 updates,
the speech Block has approximately 29.77\% soft and 22.19\% hard-call accuracy,
compared with 76.48\% for the larger Transformer. The image Block obtains
approximately 42.1\% and 30.4\%, respectively, compared with 49.9\% for the
Transformer. Their Block parameter counts are 32,280 and 40,793; the Transformer
counts are 105,379 and 110,116. These static development tests exercise the
same core through learned adapters. They do not yet test full-class continual
recognition, and the soft/hard gaps show that inference selection remains a
material architectural choice.

The central remaining obligation is continual acquisition with complete
retention and resource accounting. The continuous Block currently resets
per-example memory, recomputes observed prefixes and has no general shared-update
protection certificate. Earlier finite monodromy and binding-closure results
retain their narrower scope in the appendix. Any integration must explain
what operation or task relation is protected in the new state space and
measure its benefit against ordinary constraints. An independent new primitive
is not a prerequisite for an important empirical result; a broad reproducible
advantage still requires a frozen architecture, strong appropriate controls,
and the full evaluation rather than a favorable subset of development scores.
